% Part: first-order-logic
% Chapter: beyond
% Section: introduction

\documentclass[../../../include/open-logic-section]{subfiles}

\begin{document}

\olfileid{fol}{byd}{int}

\olsection{সারসংক্ষেপ}

প্রথম-ক্রমের যুক্তিবিদ্যাই আগ্রহের একমাত্র যৌক্তিক ব্যবস্থা নয়:
এর বহু সম্প্রসারণ ও রূপভেদ আছে। সাধারণত একটি যুক্তিবিদ্যায় থাকে কোনো
ভাষার আনুষ্ঠানিক নির্দিষ্টকরণ, সাধারণত---তবে সব সময় নয়---একটি
অবরোহী ব্যবস্থা, এবং সাধারণত---তবে সব সময় নয়---একটি অভিপ্রেত
অর্থতত্ত্ব। কিন্তু শব্দটির এই কারিগরি ব্যবহার একটি স্বাভাবিক প্রশ্ন
তোলে: স্বজ্ঞাত বা দার্শনিক অর্থে ব্যবহৃত “যুক্তি” শব্দটির সঙ্গে
প্রথম-ক্রমের যুক্তিবিদ্যা নয় এমন যুক্তিবিদ্যাগুলির সম্পর্ক কী? নিচে
বর্ণিত সব ব্যবস্থাই কোনো না কোনো ধরনের যুক্তিবিচারের মডেল হিসেবে
নির্মিত; কিন্তু কোন বৈশিষ্ট্য তাদের যৌক্তিক করে, তা কি বলা যায়?

সহজ কোনো উত্তর সামনে নেই। “যুক্তি” শব্দটি নানা অর্থে ও নানা প্রসঙ্গে
ব্যবহৃত হয় এবং “সত্য”-এর ধারণার মতো এই ধারণাটিও বহু দার্শনিক অবস্থান
থেকে বিশ্লেষিত হয়েছে। যেমন, কেউ যৌক্তিক বিচারবুদ্ধির লক্ষ্য বলতে
বুঝতে পারেন কোন বিবৃতিগুলি আবশ্যিকভাবে সত্য, অভিজ্ঞতাপূর্বভাবে সত্য,
অ-যৌক্তিক পদগুলির ব্যাখ্যা থেকে স্বাধীনভাবে সত্য, তাদের রূপের কারণেই
সত্য, অথবা ভাষাগত প্রথার বলে সত্য তা নির্ণয় করা; আর এই প্রতিটি
ধারণাই বিস্তর ব্যাখ্যা দাবি করে। কেবল গণিতে ব্যবহৃত যুক্তিবিদ্যায়
মনোযোগ সীমিত করলেও তার পরিসর নিয়ে সামান্যই ঐকমত্য আছে। যেমন,
\textit{Principia Mathematica}-তে রাসেল ও হোয়াইটহেড ফ্রেগে-প্রবর্তিত
{\em যুক্তিবাদী} ধারায় যুক্তির ভিত্তিতে গণিত গড়ে তোলার চেষ্টা
করেছিলেন। তাঁদের যৌক্তিক ব্যবস্থা ছিল নিচে বর্ণিতটির অনুরূপ এক ধরনের
উচ্চতর-টাইপের যুক্তিবিদ্যা। শেষ পর্যন্ত তাঁদের এমন কিছু স্বতঃসিদ্ধ
প্রবর্তন করতে হয়েছিল যেগুলি অধিকাংশ মানদণ্ডে বিশুদ্ধ যৌক্তিক বলে মনে
হয় না (বিশেষত অসীমতার স্বতঃসিদ্ধ ও লঘুকরণ-স্বতঃসিদ্ধ); তবু কেউ মনে
করতে পারেন যে উচ্চতর-ক্রমের কোনো কোনো যুক্তিবিচারকে যৌক্তিক বলে মেনে
নেওয়া উচিত। বিপরীতে, কুয়াইন তাঁর সত্তাতত্ত্বে “প্রস্তাবনা”-কে
আলোচনার বৈধ বস্তু হিসেবে স্বীকার করেন না এবং যুক্তি দেন যে
দ্বিতীয়-ক্রমের ও উচ্চতর-ক্রমের যুক্তিবিদ্যা আসলে ভেড়ার ছদ্মবেশে
সেটতত্ত্ব; অন্যভাবে বললে, বিধেয়ের উপর পরিমাণায়নসমৃদ্ধ ব্যবস্থাগুলি
বিশুদ্ধ যৌক্তিক নয়।

আপাতত এমন দার্শনিক প্রশ্ন বর্ষার দিনের জন্য তুলে রাখাই ভালো। নিচের
ব্যবস্থাগুলিকে শুধু বিভিন্ন ধরনের---যৌক্তিক বা অন্য---যুক্তিবিচারের
আনুষ্ঠানিক আদর্শায়ন হিসেবে ভাবা যাক।

\end{document}
